% SPDX-License-Identifier: Apache-2.0
\section{Findings}
\label{sec:findings}

The first run, on 2026-10-07 against TxHDL commit \code{894cc2e3b296},
found two departures from the specification. Both are filed as issues
against TxHDL, and both are in the privileged architecture. In that
run, every unprivileged instruction the core implements matched Sail
on every value the tests check. A later run that fails anything else
shows it in Section~\ref{sec:results} and in
Appendix~\ref{sec:failures}, as a failure that is not known.

\subsection{\code{mstatus.TVM}, \code{TW} and \code{TSR} cannot be set}

The test \code{SvSm\_sv\_mstatus\_tvm\_test-00} sets \code{mstatus.MPP}
to 3 and \code{mstatus.TVM} to 1, then reads \code{mstatus} back. Sail
reads \code{0x00101800}. Vreteno reads \code{0x00001800}: the TVM bit,
bit 20, did not stick. The test's own console output is in
Appendix~\ref{sec:failures}.

The cause is in the core's write mask for \code{mstatus},
\code{MSTATUS\_W = 0x000e\_19aa}, which leaves out bits 20, 21 and 22:
TVM, TW and TSR. The privileged specification requires all three to be
writable when supervisor mode is implemented~\cite{priv}. Without them,
machine mode cannot trap a supervisor's \code{satp} accesses and
\code{sfence.vma}, nor its \code{wfi} and \code{sret}. A hypervisor or
a monitor that emulates those instructions relies on these traps.
The issue is TxHDL issue 1347~\cite{issue1347}.

\subsection{No \code{menvcfg} and no \code{senvcfg}}

Version 1.12 of the privileged specification adds \code{menvcfg}, which
a core must have when it implements user mode, and \code{senvcfg},
which it must have when it implements supervisor mode. Vreteno has
neither, and an access to either raises an illegal instruction
exception. Described as a 1.12 core, Vreteno fails every test: the
tests write \code{menvcfg} while they boot. Described as a 1.11 core,
it passes all but one.

The core does have \code{mstatush}, which is also new in 1.12, so it
sits between the two versions. This report describes it as 1.11, and
the Svade suite, which needs \code{menvcfgh}, does not run. The issue
is TxHDL issue 1348~\cite{issue1348}. Adding the two registers, with
their fields read-only zero where the core lacks the feature, would
let the core be described as 1.12 and the Svade suite run.

\subsection{Since the first run}

TxHDL fixed both issues by commit \code{48daf9778add}, and the run of
2026-10-08 against that commit passes
\code{SvSm\_sv\_mstatus\_tvm\_test-00}. The core now has
\code{menvcfg} and \code{senvcfg} too, but this report still
describes it as a 1.11 core, so the Svade suite still does not run.

The same commit adds a data cache for the DDR3, and with it five tests
fail that passed before: the byte, halfword and word loads,
\code{I-lb-00}, \code{I-lbu-00}, \code{I-lh-00}, \code{I-lhu-00} and
\code{I-lw-00}. Each fails where it stores a new value to a line the
cache holds and loads it back at once. The load returns the value the
case before it stored. The tests run in the cached range of the
address space, and the uncached path has no such failure. The issue is
TxHDL issue 1431~\cite{issue1431}. TxHDL fixed it by commit
\code{e9681300a8b2}, and the run of 2026-10-09 against that commit
passes all five tests.

\subsection{A choice that is not a defect}

The first run had six more failures, all in tests of atomic
instructions at misaligned addresses. Vreteno raised
address-misaligned there, and the first configuration of Sail expected
an access fault. The specification allows either. The fix was in
\code{sail.json}, which now says that the core raises address-misaligned,
and the six tests pass.
